\newpage
\renewcommand{\sectioncolor}{chap6}
\section{Podsumowanie i perspektywy rozwoju}

Niniejsza praca pozwoliła na zaprojektowanie i zaimplementowanie skalowalnego systemu rezerwacji usług w modelu Mikro-SaaS. Wykorzystanie architektury mikroserwisowej opartej na platformie .NET 8 \cite{dotnet_microsoft} oraz bibliotece React 19 \cite{react_meta} umożliwiło stworzenie rozwiązania charakteryzującego się wysoką wydajnością oraz czytelnym podziałem odpowiedzialności między poszczególnymi modułami.

\subsection{Wnioski z realizacji projektu}
Proces wytwórczy oraz przeprowadzone testy automatyczne \cite{vitest_guide} pozwoliły na sformułowanie następujących wniosków inżynierskich:

\begin{itemize}
    \item \textbf{Efektywność komunikacji asynchronicznej:} Zastosowanie brokera komunikatów RabbitMQ \cite{rabbitmq_docs} pozwoliło na całkowite odseparowanie serwisu rezerwacji od modułu powiadomień. Dzięki temu proces rezerwacji terminu nie jest blokowany przez czasochłonne operacje wysyłki wiadomości e-mail.
    \item \textbf{Izolacja danych:} Model \textit{Database-per-Service} z wykorzystaniem bazy PostgreSQL \cite{postgresql_docs} zagwarantował, że błędy w obrębie jednego schematu danych (np. logów powiadomień) nie wpływają na stabilność krytycznych tabel rezerwacji.
    \item \textbf{Wydajność interfejsu:} Wykorzystanie narzędzia Vite \cite{vite_guide} oraz mechanizmów normalizacji danych wejściowych (np. telefonów i kodów pocztowych) znacząco poprawiło komfort pracy użytkownika (\textit{User Experience}) oraz jakość danych trafiających do systemu.
    \item \textbf{Bezpieczeństwo:} Implementacja mechanizmu \textit{Refresh Tokens} oraz centralna walidacja JWT na poziomie API Gateway \cite{ocelot_docs} zapewniły standard bezpieczeństwa zgodny z zaleceniami RFC 7519 \cite{jwt_rfc}.
\end{itemize}

\subsection{Kierunki dalszego rozwoju systemu}
Choć system w obecnej fazie jest w pełni funkcjonalny, architektura pozwala na jego dalszą rozbudowę w następujących obszarach:

\begin{itemize}
    \item \textbf{Integracja z kalendarzami zewnętrznymi:} Implementacja dwukierunkowej synchronizacji z Google Calendar oraz Outlook za pomocą dedykowanych API.
    \item \textbf{System płatności online:} Integracja z platformą Stripe \cite{axios_docs} w celu umożliwienia pobierania przedpłat za rezerwacje oraz zarządzania subskrypcjami firm korzystających z modelu SaaS.
    \item \textbf{Aplikacja mobilna:} Stworzenie natywnego interfejsu dla pracowników w technologii React Native, co ułatwiłoby zarządzanie harmonogramem w terenie.
    \item \textbf{Rozbudowa modułu powiadomień:} Dodanie obsługi wiadomości SMS oraz powiadomień typu \textit{Push} z wykorzystaniem technologii SignalR \cite{signalr_microsoft}.
\end{itemize}

\subsection{Podsumowanie końcowe}
Projekt udowodnił, że zastosowanie wzorców takich jak CQRS, API Gateway oraz komunikacja sterowana zdarzeniami pozwala na budowę systemów gotowych do obsługi dużego ruchu i wielu niezależnych organizacji jednocześnie. Praca nad systemem pozwoliła na zdobycie praktycznego doświadczenia w zakresie orkiestracji kontenerów Docker \cite{docker_docs} oraz automatyzacji procesów CI/CD \cite{github_actions_docs}, co stanowi solidną podstawę dla nowoczesnej inżynierii oprogramowania.
